\begin{thebibliography}{99}

\bibitem{ref1} T. A. McDonagh, M. Metra, M. Adamo, R. S. Gardner, A. Baumbach, M. B\"ohm, H. Burri, J. Butler, J. Čelutkienė, O. Chioncel \emph{et al.}, ``2021 ESC Guidelines for the diagnosis and treatment of acute and chronic heart failure,'' \emph{European Heart Journal}, vol. 42, no. 36, pp. 3599--3726, 2021.

\bibitem{ref2} P. A. Heidenreich, B. Bozkurt, D. Aguilar, L. A. Allen, J. J. Byun, M. M. Colvin, A. Deswal, M. H. Drazner, S. M. Dunlay, L. R. Evers \emph{et al.}, ``2022 AHA/ACC/HFSA Guideline for the Management of Heart Failure: A Report of the American College of Cardiology/American Heart Association Joint Committee on Clinical Practice Guidelines,'' \emph{Circulation}, vol. 145, no. 18, pp. e895--e1032, 2022.

\bibitem{ref3} Z. I. Attia, S. Kapa, F. Lopez-Jimenez, P. M. McKie, D. J. Ladewig, G. Satam, P. A. Pellikka, M. Enriquez-Sarano, P. A. Noseworthy, T. M. Munger \emph{et al.}, ``Screening for cardiac contractile dysfunction using an artificial intelligence--enabled electrocardiogram,'' \emph{Nature Medicine}, vol. 25, no. 1, pp. 70--74, 2019.

\bibitem{ref4} X. Yao, D. R. Rushlow, J. W. Inselman, R. G. McCoy, T. D. Thacher, E. M. Behnken, M. E. Bernard, S. L. Rosas, A. Akfaly, A. Misra \emph{et al.}, ``Artificial intelligence--enabled electrocardiograms for identification of patients with low ejection fraction: a pragmatic, randomized clinical trial,'' \emph{Nature} \emph{Medicine}, vol. 27, no. 5, pp. 815--819, 2021.

\bibitem{ref5} R. M. Lang, L. P. Badano, V. Mor-Avi, J. Afilalo, A. Armstrong, L. Ernande, F. A. Flachskampf, E. Foster, S. A. Goldstein, T. Kuznetsova \emph{et al.}, ``Recommendations for Cardiac Chamber Quantification by Echocardiography in Adults: An Update from the American Society of Echocardiography and the European Association of Cardiovascular Imaging,'' \emph{Journal of the American Society of} \emph{Echocardiography}, vol. 28, no. 1, pp. 1--39.e14, 2015.

\bibitem{ref6} P. Thavendiranathan, A. D. Grant, T. Negishi, J. C. Plana, Z. B. Popović, and T. H. Marwick, ``Reproducibility of Echocardiographic Techniques for Sequential Assessment of Left Ventricular Ejection Fraction and Volumes: Application to Patients Undergoing Cancer Chemotherapy,'' \emph{Journal of the} \emph{American College of Cardiology}, vol. 61, no. 1, pp. 77--84, 2013.

\bibitem{ref7} D. G. Strauss, R. H. Selvester, and G. S. Wagner, ``Defining Left Bundle Branch Block in the Era of Cardiac Resynchronization Therapy,'' \emph{The American Journal of Cardiology}, vol. 107, no. 6, pp. 927--934, 2011.

\bibitem{ref8} M. K. Das, B. Khan, S. Jacob, A. Kumar, and J. Mahenthiran, ``Significance of a Fragmented QRS Complex Versus a Q Wave in Patients With Coronary Artery Disease,'' \emph{Circulation}, vol. 113, no. 21, pp. 2495--2501, 2006.

\bibitem{ref9} I. Kardys, J. A. Kors, I. M. van der Meer, A. Hofman, D. A. M. van der Kuip, and J. C. M. Witteman, ``Spatial QRS-T angle predicts cardiac death in a general population,'' \emph{European Heart Journal}, vol. 24, no. 14, pp. 1357--1364, 2003.

\bibitem{ref10} B. Acar, G. Yi, K. Hnatkova, and M. Malik, ``Spatial, temporal and wavefront direction characteristics of 12-lead T-wave morphology,'' \emph{Medical \& Biological Engineering \& Computing}, vol. 37, no. 5, pp. 574--584, 1999.

\bibitem{ref11} R. D. Berger, E. K. Kasper, K. L. Baughman, E. Marban, H. Calkins, and G. F. Tomaselli, ``Beatto-Beat QT Interval Variability: Novel Evidence for Repolarization Lability in Ischemic and Nonischemic Dilated Cardiomyopathy,'' \emph{Circulation}, vol. 96, no. 5, pp. 1557--1565, 1997.

\bibitem{ref12} A. Vaid, K. W. Johnson, M. A. Badgeley, S. S. Somani, M. Bicak, I. Landi, A. Russak, S. Zhao, M. A. Levin, R. S. Freeman \emph{et al.}, ``Using Deep-Learning Algorithms to Simultaneously Identify Right and Left Ventricular Dysfunction From the Electrocardiogram,'' \emph{JACC: Cardiovascular Imaging}, vol. 15, no. 3, pp. 395--410, 2022.

\bibitem{ref13} V. Sangha, A. A. Nargesi, L. S. Dhingra, A. Khunte, B. J. Mortazavi, A. H. Ribeiro, E. Banina, O. Adeola, N. Garg, C. A. Brandt \emph{et al.}, ``Detection of Left Ventricular Systolic Dysfunction From Electrocardiographic Images,'' \emph{Circulation}, vol. 148, no. 9, pp. 765--777, 2023.

\bibitem{ref14} K. K. Shivashankara, Deepanshi, A. Mehri Shervedani, G. D. Clifford, M. A. Reyna, and R. Sameni, ``ECG-Image-Kit: a synthetic image generation toolbox to facilitate deep learning-based electrocardiogram digitization,'' \emph{Physiological Measurement}, vol. 45, no. 5, p. 055019, 2024.

\bibitem{ref15} H. Wu, K. H. K. Patel, X. Li, B. Zhang, C. Galazis, N. Bajaj, A. Sau, X. Shi, L. Sun, Y. Tao \emph{et al.}, ``A fully-automated paper ECG digitisation algorithm using deep learning,'' \emph{Scientific Reports}, vol. 12, no. 1, p. 20963, 2022.

\bibitem{ref16} M. A. Reyna, Deepanshi, J. T. Weigle, Z. Koscova, K. S. Campbell, S. Seyedi, A. Elola, A. Bahrami Rad, A. J. Shah, N. K. Bhatia \emph{et al.}, ``Digitization and Classification of ECG Images: The George B. Moody PhysioNet Challenge 2024,'' in \emph{Computing in Cardiology (CinC 2024)}, vol. 51, 2024.

\bibitem{ref17} A. E. Ulloa-Cerna, L. Jing, J. M. Pfeifer, S. Raghunath, J. A. Ruhl, D. B. Rocha, J. B. Leader, N. Zimmerman, G. Lee, S. R. Steinhubl \emph{et al.}, ``rECHOmmend: An ECG-Based Machine Learning Approach for Identifying Patients at Increased Risk of Undiagnosed Structural Heart Disease Detectable by Echocardiography,'' \emph{Circulation}, vol. 146, no. 1, pp. 36--47, 2022.

\bibitem{ref18} Z. I. Attia, D. M. Harmon, J. Dugan, L. Manka, F. Lopez-Jimenez, A. Lerman, K. C. Siontis, P. A. Noseworthy, X. Yao, E. W. Klavetter \emph{et al.}, ``Prospective evaluation of smartwatch-enabled detection of left ventricular dysfunction,'' \emph{Nature Medicine}, vol. 28, no. 12, pp. 2497--2503, 2022.

\bibitem{ref19} P. Bachtiger, C. F. Petri, F. E. Scott, S. Ri Park, M. A. Kelshiker, H. K. Sahemey, B. Dumea, R. Alquero, P. S. Padam, I. R. Hatrick \emph{et al.}, ``Point-of-care screening for heart failure with reduced ejection fraction using artificial intelligence during ECG-enabled stethoscope examination in London, UK: a prospective, observational, multicentre study,'' \emph{The Lancet Digital Health}, vol. 4, no. 2, pp. e117--e125, 2022.

\bibitem{ref20} A. Khunte, V. Sangha, E. K. Oikonomou, L. S. Dhingra, A. Aminorroaya, B. J. Mortazavi, A. Coppi, C. A. Brandt, H. M. Krumholz, and R. Khera, ``Detection of left ventricular systolic dysfunction from single-lead electrocardiography adapted for portable and wearable devices,'' \emph{npj Digital} \emph{Medicine}, vol. 6, no. 1, p. 124, 2023.

\bibitem{ref21} J. Li, A. D. Aguirre, V. Moura, J. Jin, C. Liu, L. Zhong, C. Sun, G. Clifford, M. B. Westover, and S. Hong, ``An Electrocardiogram Foundation Model Built on over 10 Million Recordings,'' \emph{NEJM} \emph{AI}, vol. 2, no. 7, p. AIoa2401033, 2025.

\bibitem{ref22} K. McKeen, S. Masood, A. Toma, B. Rubin, and B. Wang, ``ECG-FM: an open electrocardiogram foundation model,'' \emph{JAMIA Open}, vol. 8, no. 5, p. ooaf122, 2025.

\bibitem{ref23} T. J. Poterucha, L. Jing, R. P. Ricart, M. Adjei-Mosi, J. Finer, D. Hartzel, C. Kelsey, A. Long, D. Rocha, J. A. Ruhl \emph{et al.}, ``Detecting structural heart disease from electrocardiograms using AI,'' \emph{Nature}, vol. 644, no. 8075, pp. 221--230, 2025.

\bibitem{ref24} J. W. Hughes, L. Jing, J. Finer, D. Hartzel, C. Kelsey, A. Long, D. Rocha, J. Ruhl, T. Poterucha, and P. Elias, ``EchoNext-Mini: A Dataset and Baseline AI Model for Detecting Structural Heart Disease from Electrocardiograms,'' \emph{NEJM AI}, vol. 3, no. 5, p. AIdbp2500516, 2026.

\bibitem{ref25} A. Y. Hannun, P. Rajpurkar, M. Haghpanahi, G. H. Tison, C. Bourn, M. P. Turakhia, and A. Y. Ng, ``Cardiologist-level arrhythmia detection and classification in ambulatory electrocardiograms using a deep neural network,'' \emph{Nature Medicine}, vol. 25, no. 1, pp. 65--69, 2019.

\bibitem{ref26} A. H. Ribeiro, M. H. Ribeiro, G. M. M. Paixão, D. M. Oliveira, P. R. Gomes, J. A. Canazart, M. P. S. Ferreira, C. R. Andersson, P. W. Macfarlane, W. Meira \emph{et al.}, ``Automatic diagnosis of the 12-lead ECG using a deep neural network,'' \emph{Nature Communications}, vol. 11, no. 1, p. 1760, 2020.

\bibitem{ref27} K. He, X. Zhang, S. Ren, and J. Sun, ``Deep Residual Learning for Image Recognition,'' in \emph{Pro-} \emph{ceedings of the IEEE Conference on Computer Vision and Pattern Recognition (CVPR)}, 2016, pp. 770--778.

\bibitem{ref28} M. Tan and Q. Le, ``EfficientNet: Rethinking Model Scaling for Convolutional Neural Networks,'' in \emph{Proceedings of the 36th International Conference on Machine Learning (ICML)}, ser. Proceedings of Machine Learning Research, vol. 97. PMLR, 2019, pp. 6105--6114. [Online]. Available: \url{https://proceedings.mlr.press/v97/tan19a.html}

\bibitem{ref29} A. Howard, M. Sandler, B. Chen, W. Wang, L.-C. Chen, M. Tan, G. Chu, V. Vasudevan, Y. Zhu, R. Pang \emph{et al.}, ``Searching for MobileNetV3,'' in \emph{Proceedings of the IEEE/CVF International Con-} \emph{ference on Computer Vision (ICCV)}, 2019, pp. 1314--1324.

\bibitem{ref30} Z. Liu, H. Mao, C.-Y. Wu, C. Feichtenhofer, T. Darrell, and S. Xie, ``A ConvNet for the 2020s,'' in \emph{Proceedings of the IEEE/CVF Conference on Computer Vision and Pattern Recognition (CVPR)}, 2022, pp. 11 966--11 976.

\bibitem{ref31} S. Hochreiter and J. Schmidhuber, ``Long Short-Term Memory,'' \emph{Neural Computation}, vol. 9, no. 8, pp. 1735--1780, 1997.

\bibitem{ref32} Z. I. Attia, S. Kapa, X. Yao, F. Lopez-Jimenez, T. L. Mohan, P. A. Pellikka, R. E. Carter, N. D. Shah, P. A. Friedman, and P. A. Noseworthy, ``Prospective validation of a deep learning electrocardiogram algorithm for the detection of left ventricular systolic dysfunction,'' \emph{Journal of Cardiovascular Elec-} \emph{trophysiology}, vol. 30, no. 5, pp. 668--674, 2019.

\bibitem{ref33} D. Adedinsewo, R. E. Carter, Z. Attia, P. Johnson, A. H. Kashou, J. L. Dugan, M. Albus, J. M. Sheele, F. Bellolio, P. A. Friedman \emph{et al.}, ``Artificial Intelligence-Enabled ECG Algorithm to Identify Patients With Left Ventricular Systolic Dysfunction Presenting to the Emergency Department With Dyspnea,'' \emph{Circulation: Arrhythmia and Electrophysiology}, vol. 13, no. 8, p. e008437, 2020.

\bibitem{ref34} J.-M. Kwon, K.-H. Kim, K.-H. Jeon, H. M. Kim, M. J. Kim, S.-M. Lim, P. S. Song, J. Park, R. K. Choi, and B.-H. Oh, ``Development and Validation of Deep-Learning Algorithm for Electrocardiography-Based Heart Failure Identification,'' \emph{Korean Circulation Journal}, vol. 49, no. 7, pp. 629--639, 2019.

\bibitem{ref35} V. Sangha, B. J. Mortazavi, A. D. Haimovich, A. H. Ribeiro, C. A. Brandt, D. L. Jacoby, W. L. Schulz, H. M. Krumholz, A. L. P. Ribeiro, and R. Khera, ``Automated multilabel diagnosis on electrocardiographic images and signals,'' \emph{Nature Communications}, vol. 13, no. 1, p. 1583, 2022.

\bibitem{ref36} P. Elias, T. J. Poterucha, V. Rajaram, L. M. Moller, V. Rodriguez, S. Bhave, R. T. Hahn, G. Tison, S. A. Abreau, J. Barrios \emph{et al.}, ``Deep Learning Electrocardiographic Analysis for Detection of Left-Sided Valvular Heart Disease,'' \emph{Journal of the American College of Cardiology}, vol. 80, no. 6, pp. 613--626, 2022.

\bibitem{ref37} S. Varma and R. Simon, ``Bias in error estimation when using cross-validation for model selection,'' \emph{BMC Bioinformatics}, vol. 7, no. 1, p. 91, 2006.

\bibitem{ref38} C. Nadeau and Y. Bengio, ``Inference for the Generalization Error,'' \emph{Machine Learning}, vol. 52, no. 3, pp. 239--281, 2003.

\bibitem{ref39} S. Holm, ``A Simple Sequentially Rejective Multiple Test Procedure,'' \emph{Scandinavian Journal of} \emph{Statistics}, vol. 6, no. 2, pp. 65--70, 1979. [Online]. Available: \url{https://www.jstor.org/stable/4615733}

\end{thebibliography}
